% VI36-DETAILED-PROBLEM-01
\begin{problem}[VI.36.P1: Reconstruct \(\mathbb P^1\) by gluing]\label{prob:vi36-01}
Starting from two copies \(X_0=\Spec k[t]\) and \(X_1=\Spec k[u]\), glue \(D(t)\) to \(D(u)\) by \(u=t^{-1}\). Prove that the resulting scheme is isomorphic to \(\mathbb P^1_k\).
\end{problem}

\ifdefined\FullProblemDossiers
\begin{solution}
Let
\[
S=k[x_0,x_1].
\]
The standard cover of
\[
\mathbb P^1_k=\Proj S
\]
is
\[
U_0=D_+(x_0),
\qquad
U_1=D_+(x_1).
\]
VI/35 gives
\[
U_0\cong\Spec S_{(x_0)}
=
\Spec k[x_1/x_0].
\]
Set
\[
t=x_1/x_0.
\]
Likewise
\[
U_1\cong\Spec k[x_0/x_1],
\qquad
u=x_0/x_1.
\]

Their intersection is
\[
D_+(x_0x_1).
\]
On the first chart this is \(D(t)\), with ring \(k[t,t^{-1}]\), and on the second it is \(D(u)\), with ring
\(k[u,u^{-1}]\). The homogeneous ratios satisfy
\[
u=t^{-1}.
\]
Thus the standard Proj cover has exactly the same affine pieces and overlap identification as the proposed gluing.
By uniqueness in the scheme-gluing theorem,
\[
\boxed{
X_0\cup_{u=t^{-1}}X_1
\cong
\mathbb P^1_k.
}
\]
\end{solution}
\fi
